% ============================================================
% 02-neuravnotezenost-i-smote.tex — PRVA VERZIJA
% ============================================================

\subsection{Klasifikacija i problem neuravnote\v{z}enosti klasa}

U ovom potpoglavlju bit \'{c}e opisano \v{s}to predstavlja klasifikacija
\cite{duda2000,theodoridis2008} i \v{s}to ima za cilj, ograni\v{c}avaju\'{c}i
se pritom na skupove podataka $\mathcal{A} \subset \mathbb{R}^m$ gdje su
sve zna\v{c}ajke koje opisuju podatke realne. Spomenut \'{c}e se neke
uobi\v{c}ajene primjene klasifikacije u razli\v{c}itim podru\v{c}jima,
uklju\v{c}uju\'{c}i medicinsku dijagnostiku \cite{odajima2014,kamel2019} i
druge domene.

Zatim se opisuje \v{s}to je neuravnote\v{z}enost klasa i na koji
na\v{c}in utje\v{c}e na problem klasifikacije --- u pravilu, rezultira
pristrano\v{s}\'{c}u klasifikatora prema ve\'{c}inskoj klasi
\cite{he2009,lopez2013}. Navest \'{c}e se za\v{s}to se radi o zna\v{c}ajnom
problemu i za\v{s}to je nu\v{z}no rukovati njime na odgovaraju\'{c}i
na\v{c}in, posebice jer je manjinska klasa \v{c}esto upravo ona od
interesa u stvarnim primjenama. Formalno \'{c}e se definirati omjer
neuravnote\v{z}enosti $IR = N_{\text{maj}} / N_{\text{min}}$ te
kategorizacija skupova prema vrijednosti IR-a.

Ukratko se opisuju poduzorkovanje i preuzorkovanje kao na\v{c}ini
rukovanja navedenim problemom. Iako nema univerzalno boljeg pristupa i
u\v{c}inkovitost ovisi o skupu podataka i klasifikatoru, preuzorkovanje
se \v{c}e\v{s}\'{c}e koristi jer kod poduzorkovanja postoji opasnost od
gubljenja bitnih uzoraka ve\'{c}inske klase \cite{bajer2019}. U ovom
dijelu ukomponirat \'{c}e se i podjela skupa podataka na podskup za
treniranje i testiranje.

U odnosu na postoje\'{c}u literaturu, dodatno \'{c}e se obratiti
pozornost na ote\v{z}avaju\'{c}e faktore koji prate problem
neuravnote\v{z}enosti: preklapanje klasa, unutar-klasni disbalans,
prisutnost \v{s}uma u podacima, mali apsolutni broj manjinskih primjera
te manifold struktura podataka \cite{krawczyk2016}.

\subsubsection{Neki popularni algoritmi za klasifikaciju}

Postoje brojni algoritmi za klasifikaciju
\cite{duda2000,theodoridis2008} --- nisu svi jednako prikladni za sve
probleme i odabir klasifikatora (engl.\ model selection) predstavlja
zaseban problem u strojnom u\v{c}enju. U nastavku se navodi osam
klasifikatora koji \'{c}e se koristiti u eksperimentalnom dijelu
rada, uz princip rada, klju\v{c}ne parametre te prednosti i nedostatke:

\begin{itemize}
    \item Stablo odluke (engl.\ Decision Tree, DT) --- rekurzivno dijeli
    podatke temeljem vrijednosti zna\v{c}ajki, sklono prekomjernom
    prilago\dj{}avanju.
    \item Slu\v{c}ajna \v{s}uma (engl.\ Random Forest, RF) --- ansambl
    stabala odluke s baggingom, smanjuje varijancu i robusnija je od
    pojedina\v{c}nog stabla.
    \item XGBoost (engl.\ Extreme Gradient Boosting) --- gradient boosting
    ansambl s regularizacijom, \v{c}esto posti\v{z}e vrhunske rezultate
    na tabli\v{c}nim podacima; ima ugra\dj{}enu podr\v{s}ku za
    neuravnote\v{z}ene klase putem parametra \texttt{scale\_pos\_weight}.
    \item Logisti\v{c}ka regresija (LR) --- linearni model s vjerojatnosnim
    izlazom, jednostavan i \v{c}esto dobar baseline.
    \item Stroj s potpornim vektorima (engl.\ Support Vector Machine, SVM)
    --- maksimizira marginu izme\dj{}u klasa, RBF jezgra omogu\'{c}uje
    nelinearnu separaciju.
    \item Algoritam $k$-najbli\v{z}ih susjeda ($k$-NN) --- klasificira
    temeljem klasa $k$ najbli\v{z}ih primjera, osjetljiv na distribuciju
    i skalu zna\v{c}ajki.
    \item Gaussov naivni Bayes (GNB) --- pretpostavlja normalnu distribuciju
    zna\v{c}ajki i me\dj{}usobnu nezavisnost, brz i jednostavan.
    \item Vi\v{s}eslojni perceptron (engl.\ Multi-Layer Perceptron, MLP)
    --- jednostavna neuronska mre\v{z}a s jednim skrivenim slojem,
    predstavlja osnovnu deep learning referentnu to\v{c}ku u usporedbi.
\end{itemize}

\subsubsection{Vrednovanje u\v{c}inkovitosti algoritama klasifikacije}

Objasnit \'{c}e se \v{s}to predstavlja matrica zbunjenosti (engl.\
confusion matrix) i kako se dobiva. Uobi\v{c}ajeno se u\v{c}inkovitost
klasifikatora vrednuje to\v{c}no\v{s}\'{c}u klasifikacije (engl.\
classification accuracy), no ona nije prikladna kod neuravnote\v{z}enih
klasa --- klasifikator mo\v{z}e posti\'{c}i visoku to\v{c}nost
predvi\dj{}aju\'{c}i isklju\v{c}ivo ve\'{c}insku klasu.

Zatim \'{c}e se opisati prikladnije mjere za neuravnote\v{z}ene podatke,
polaze\'{c}i od elemenata matrice zbunjenosti (TP, TN, FP, FN):
F-1 mjera kao harmonijska sredina preciznosti i odziva, G-Mean
kao geometrijska sredina odziva i specifi\v{c}nosti, AUC-ROC kao povr\v{s}ina
ispod ROC krivulje neovisna o pragu odluke, Balanced Accuracy kao
aritmeti\v{c}ka sredina odziva i specifi\v{c}nosti te Matthewsov
koeficijent korelacije (MCC) koji uzima u obzir sva \v{c}etiri elementa
matrice zbunjenosti. Za svaku mjeru dat \'{c}e se formula, interpretacija
i rasprava o prikladnosti. Ograni\v{c}it \'{c}e se na slu\v{c}aj binarne
klasifikacije gdje je pozitivna klasa uobi\v{c}ajeno manjinska.
Pritom \'{c}e se koristiti literatura \cite{sokolova2009,japkowicz2011,lopez2013,he2009}.

\subsection{Pregled pristupa rje\v{s}avanju problema}

U ovom potpoglavlju dat \'{c}e se krovni pregled \v{c}etiriju kategorija
pristupa rje\v{s}avanju problema neuravnote\v{z}enosti klasa, pri \v{c}emu
\'{c}e se detaljni opisi pojedinih metoda obraditi u tre\'{c}em poglavlju
\cite{fernandez2018book}.

Prvo, metode na razini podataka koje obuhva\'{c}aju
preuzorkovanje manjinske klase, poduzorkovanje ve\'{c}inske klase i
njihove hibridne kombinacije. Prednost ovih metoda je neovisnost o
klasifikatoru i jednostavnost primjene \cite{bajer2019,dudjak2020}.
Drugo, metode na razini algoritma koje uklju\v{c}uju cost-sensitive
u\v{c}enje i modifikacije klasifikatora \cite{elkan2001}. Tre\'{c}e,
ansambl metode koje kombiniraju resampliranje s bagging i boosting
pristupima \cite{galar2012}. \v{C}etvrto, deep learning pristupi koji
obuhva\'{c}aju modifikacije loss funkcija \cite{lin2017} i generativne
modele \cite{goodfellow2014}.

\subsection{Algoritam SMOTE}

U ovom potpoglavlju detaljno \'{c}e se opisati temeljni SMOTE algoritam
\cite{chawla2002}. Prvo \'{c}e se ukratko opisati nasumi\v{c}no
preuzorkovanje kao motivacija za nastanak SMOTE-a --- kod nasumi\v{c}nog
preuzorkovanja dupliciraju se postoje\'{c}i primjeri manjinske klase,
\v{s}to \v{c}esto dovodi do prekomjernog prilago\dj{}avanja. SMOTE ovaj
problem rje\v{s}ava generiranjem novih, sinteti\v{c}kih primjera.

Zatim \'{c}e se dati detaljan opis koraka algoritma: za svaki primjer
manjinske klase $x_i$ pronalazi se $k$ najbli\v{z}ih susjeda unutar iste
klase (Euklidskom udaljeno\v{s}\'{c}u); nasumi\v{c}no se odabire jedan
susjed $x_{nn}$; novi sinteti\v{c}ki primjer generira se linearnom
interpolacijom prema formuli $x_{new} = x_i + \lambda \cdot (x_{nn} - x_i)$,
gdje je $\lambda \sim U(0,1)$. Postupak se ponavlja dok se ne postigne
\v{z}eljeni broj sinteti\v{c}kih primjera. Bit \'{c}e prikazan i
pseudo-k\^{o}d algoritma.

Spomenut \'{c}e se parametri algoritma: broj najbli\v{z}ih susjeda $k$
(uobi\v{c}ajeno $k = 5$ prema preporuci iz izvornog rada) i postotak
preuzorkovanja $N$ (naj\v{c}e\v{s}\'{c}e se cilja na potpuno balansiranje).
Naglasit \'{c}e se da izvorni SMOTE generira isklju\v{c}ivo to\v{c}ke na
linijskim segmentima, \v{s}to ograni\v{c}ava raznolikost sinteti\v{c}kih
uzoraka i motivira razvoj geometrijskih pro\v{s}irenja opisanih u
tre\'{c}em poglavlju.

\subsubsection{Prednosti i nedostaci}

Navest \'{c}e se i kratko obrazlo\v{z}iti prednosti i nedostaci algoritma
SMOTE. Me\dj{}u prednostima istaknut \'{c}e se: smanjenje overfittinga u
odnosu na nasumi\v{c}no preuzorkovanje, pro\v{s}irenje regije odluke
manjinske klase, jednostavnost implementacije, neovisnost o klasifikatoru
te fleksibilnost u pogledu omjera preuzorkovanja.

Me\dj{}u nedostacima obradit \'{c}e se: prekomjerna generalizacija jer
algoritam generira primjere neovisno o distribuciji ve\'{c}inske klase
\cite{stefanowski2007}, osjetljivost na \v{s}um i outlier-e, fiksni broj
susjeda $k$ koji sve primjere tretira jednako, problem s manifold podacima
\cite{bellinger2016}, pote\v{s}ko\'{c}e s opravdavanjem sinteti\v{c}kih
primjera u domenama poput medicine te smanjena u\v{c}inkovitost u
visokodimenzionalnim prostorima. Pregled navedenih ograni\v{c}enja
temeljit \'{c}e se na radu \cite{fernandez2018}.
